The pipeline achieves its primary recall target on both the CREMI and XPRESS
benchmarks, but several open research directions remain. The most important
are outlined below.

\section{Precision Improvement}

Addressing the precision gap (currently $\approx$0.003 on XPRESS) is the
highest-priority research problem. The five remaining false negatives on the
training volume all fail the CurvatureRule at gaps of 66--987~nm and have
composite scores of 0.41--0.64.  Because the direction estimate at short range
is sensitive to skeleton geometry near the gap boundary, these cases appear
unrecoverable without access to raw image data to disambiguate curvature.
Three approaches are under consideration for the broader precision challenge:

\begin{enumerate}
  \item \textbf{Gap-dependent accept thresholds at recall $\leq 0.90$.}
        The per-bin ceiling analysis (Section~\ref{sec:per-bin-threshold})
        quantifies the ceiling: at recall 0.90, per-bin thresholds can
        reach precision 0.005 (+51\%); at recall 0.85, 0.007 (+104\%);
        at recall 0.70, 0.017 (+421\%).  These are concrete deployment
        levers for use cases where trading some recall for precision is
        acceptable.  At recall 0.99, however, the ceiling is only $+$1.3\%
        --- so a per-bin policy is \emph{not} the right tool at the pipeline's
        current operating point.

  \item \textbf{Features orthogonal to (composite, gap\_bin).}
        The per-bin ceiling of 0.003 at recall 0.99 means no amount
        of calibration or re-thresholding on the current composite score
        can materially improve high-recall precision.  Only features
        outside the \((\text{composite}, \text{gap\_bin})\) subspace can
        break this ceiling.  Two empirical data points show this is
        possible: Experiment~20's per-fragment degree reached 0.037 at
        recall 0.85 (5$\times$ the per-bin ceiling of 0.007 at
        the same recall), and Section~\ref{sec:density-reranker}'s
        neighborhood density sits on the ceiling at recall 0.85.  A
        multi-feature gap-stratified ensemble combining degree, density,
        composite sub-scores, and per-bin calibration is therefore the
        most promising single path to precision gains at recall $\geq 0.99$.
        Additional orthogonal cues worth exploring include myelin sheath
        consistency, local radius variation, and branch topology in a
        small window around the gap.

  \item \textbf{Gap-stratified ML filter retraining.} The Experiment~20
        classifier was trained on degree features over the full gap range
        and is effective at recall 0.85 but not at 0.99.  Because the
        feature distribution and operating point differ sharply between
        short- and long-range bins (Section~\ref{sec:gap-stratified}),
        training a separate classifier per bin --- or a single model
        conditioned explicitly on gap distance --- should improve the
        precision ceiling at high recall.

  \item \textbf{Agglomerative confidence scoring.} Rather than accepting all
        candidates above a fixed composite threshold, a global agglomerative
        approach could select the best non-conflicting set of merges across the
        entire volume, trading some recall for large precision gains.
\end{enumerate}

\section{Coverage Expansion}

The current 247 unreachable ground-truth pairs fall into two categories: 172 with
true gap $> 2000$~nm, and 75 involving PCA-only endpoint fragments whose
bounding-box tip-to-tip distance does not span the gap. Two options:

\begin{itemize}
  \item Raise \texttt{max\_endpoint\_search\_nm} from 2000 to 3000~nm.
        Experiment~13 showed this causes recall regression, but that was before
        distance-conditioned scoring (Exp~14) and CurvatureRule skip (Exp~16)
        were added. A re-test at 3000~nm with the current scoring may give
        different results.

  \item Raise \texttt{max\_skeleton\_voxels} to cover the 63 PCA-only fragments,
        each requiring 5--30 seconds of TEASAR time ($\sim 13$ minutes additional
        runtime total). Better skeleton coverage would improve endpoint accuracy
        for large fragments.
\end{itemize}

\section{Raw-Image-Aware Curvature Disambiguation}

All five remaining CurvatureRule false negatives have strong alignment
($0.40$--$0.66$) and composite scores well above the acceptance threshold, but
fail a geometric consistency check that cannot be resolved from skeleton data
alone.  A targeted post-validation pass that consults the raw EM or XNH volume
in a small window around the gap --- checking for membrane continuity, myelin
sheath persistence, or boundary contrast --- could rescue these cases without
introducing new false positives elsewhere in the pipeline.  This would be the
first voxel-level step in an otherwise graph-only pipeline and is therefore
best treated as an optional downstream module rather than a core stage.

\section{Extending Beyond Tubular Morphologies}

The pipeline's geometric scoring assumes consistent local tangent alignment
between the two fragments of a split pair, which holds for myelinated axons
but not for orthogonal-break structures such as orphaned dendritic spines
in cortical gray matter (a regime in which higher-resolution EM datasets
commonly operate). Adapting the pipeline to these morphologies would
require one of: (a) a morphology-aware scoring module that estimates
structure type (axon / dendrite / spine) from local skeleton statistics
and selects a corresponding weight vector per candidate; (b) a learned
merge-proposal classifier trained on orthogonal-break examples from
proofread EM data (e.g., MICrONS~\cite{microns2025}) that replaces the
alignment/curvature rules for fragments flagged as non-tubular; or (c) a
raw-image-conditioned pass (similar in spirit to
Section~\ref{sec:per-bin-threshold}'s futurework entry on curvature
disambiguation) that relies on local boundary continuity rather than
skeleton direction for orthogonal breaks. Option (b) is the most
promising because proofread gray-matter training data is now available at
scale.

\section{Merge-Error Detection}

The current pipeline is scoped to split-correction only; it does not
propose cuts to resolve merge errors (two distinct neurons fused into
one segment by the upstream segmenter). Merge errors are the more serious
failure mode in downstream connectivity analysis because they fabricate
synapses between genuinely unconnected neurons. A natural extension is a
symmetric cut-proposal module: for each fragment, identify candidate cut
locations using skeleton branch-point topology, local radius changes, or
learned morphological cues, and score each cut with a three-outcome
validation that mirrors the current merge validation. This would make the
pipeline a complete split-plus-merge proofreader. Recent work such as
NEURD~\cite{celii2025neurd} and ConnectomeBench~\cite{brown2025connectomebench}
addresses merge detection from the mesh/LLM side; a geometric-rules
approach consistent with this pipeline's design philosophy is open.

\section{XPRESS Challenge Submission}

The pipeline results have not yet been submitted to the XPRESS challenge portal
(\url{https://xpress.grand-challenge.org/}) for official evaluation. A formal
submission would provide an external, blind evaluation and place the results
in the context of other participating systems. The held-out validation recall
of 0.9879 suggests the configuration is ready for such a submission without
further tuning.

\section{Scaling to Full Cortex Volumes}

The XPRESS training and validation volumes ($699^3$ voxels each) are at the
small end of modern connectomics datasets; a full-cortex reconstruction contains
tens to hundreds of teravoxels.  The pipeline's current chunked fragment
extraction and KD-tree graph construction scale linearly in volume size, but
memory pressure and I/O patterns at that scale have not been characterized.
A natural continuation of this work is a benchmarking study on a larger
subvolume, together with a distributed variant of the skeleton-node graph
construction that keeps only fragment-local skeleton summaries resident in
memory.
